\documentclass[10pt,a4paper]{amsart}
\usepackage[margin=19mm]{geometry}
\usepackage{amsmath,amssymb,mathtools}
\usepackage[scr=rsfs,cal=euler]{mathalfa}
\usepackage{xfrac}
\usepackage[unicode,hidelinks,bookmarks=false]{hyperref}
\newcommand{\ie}{i.e.\ }
\newcommand{\cf}{cf.\ }
\newcommand{\resp}{respectively\ }
\newcommand{\Subsection}{\subsection}
\providecommand{\nobreakdash}{\nobreak\hbox{-}}
\setlength{\emergencystretch}{2em}
\multlinegap0pt
\usepackage{amssymb}
\newtheorem{lem}{Lemma}
\renewcommand{\le}{\leqslant}\renewcommand{\leq}{\leqslant}
\title{Boundary value problems on $p$-adic analytic manifolds}
\author{Patrick Erik Bradley}
\date{Source: arXiv:2605.16590v1 (CC BY 4.0)}
\begin{document}
\maketitle

\noindent\emph{Source and licence.} Boundary value problems on $p$-adic analytic manifolds. Version arXiv:2605.16590v1 (CC BY 4.0), distributed by arXiv under the Creative Commons Attribution 4.0 International licence, http://creativecommons.org/licenses/by/4.0/, as recorded in the arXiv licence metadata for this exact version and shown on the abstract page https://arxiv.org/abs/2605.16590v1. Full licence notice: \texttt{licenses/arxiv-2605.16590-CC-BY-4.0.txt}.\par\smallskip

\noindent\textit{Editorial scope.} Counted unit: the article's lem lem (source0.tex lines 496--498). The article's complete original proof of it is delivered with it (source0.tex lines 500--523, one \texttt{proof} environment of 24 lines). No mathematics is written, repaired or re-derived: the statement and the proof are the article's own lines, in the article's own notation, and no retained line is substituted or re-flowed.  Retained uncounted as the source-local context the proof needs: lines 486--489.  Nothing was omitted from any retained span, so the package carries no omission record.  No retained line contains a citation, so the wrapper carries no bibliography and no bibliography entry.  No retained line contains a figure environment, an image-inclusion command or drawing code, so no image is delivered and the package declares no asset dependency.  Editorial formatting only, in addition: an AMS \texttt{amsart} preamble that restates the article's own declarations and macros used by the retained lines, namely the environments \texttt{lem} on separate counters, together with the standard packages those lines need.  The retained environments print in this document as Lemma 1 in order of appearance, whereas the article prints the same items with the numbering of its own full document.  The complete native source of the article as distributed in the arXiv e-print for this exact version is bundled unchanged as \texttt{sources/200726/source0.tex}.
\begin{align}\label{pathlength}
\ell_\Lambda(\gamma)=\sum\limits_{\mathfrak{s}\in\gamma_x}\mu_\Lambda(U_{\mathfrak{s}}^\circ(x))+\sum\limits_{\mathfrak{s}\in\gamma_y}\mu_\Lambda(U_{\mathfrak{s}}^\circ(x))
+\sum\limits_{\sigma\in\gamma_{xy}^\circ}\mu_\Lambda(U_\sigma)\,.
\end{align}

\begin{lem}
The geodetic distance $d_\Lambda(x,y)$ extends to a metric on $X$.
\end{lem}

\begin{proof}
The extension is evidently to be taken as $d_\Lambda(x,x)=0$ for $x\in X$. If $x\neq y$, then any path $\gamma\colon x\leadsto y$ contains an element of $\mathfrak{B}(X,\mathcal{A})$. Since the nerve complex is finite, and its complement in $\mathfrak{B}(X,\mathcal{A})$ is a disjoint union of finitely many rooted trees, it follows that there are only finitely many paths between $x$ and $y$, and each has finite length. Hence, the infimum is a positive number, and thus $(x,y)\mapsto d_\Lambda(x,y)$ is a positive definite function.  

\smallskip
The symmetry of $d_\Lambda(x,y)$ is immediate.

\smallskip
The triangle inequality follows thus: first,
\[
d_\Lambda(x,y)=\ell_\Lambda(\gamma_x)+d_\Lambda(\sigma(x),\sigma(y))+\ell_\Lambda(\gamma_y)\,,
\]
where $\gamma_x$ is the geodesic in the rooted tree of balls containing $x$ with root $\sigma(x)$ in the nerve complex $N(\mathcal{A})$. Likewise, $\gamma_y$ is the geodesic connecting simplex $\sigma(y)$ with $y$. Further, $d_\Lambda(\sigma(x),\sigma(y))$ is the third summand in the definition of $\ell_\Lambda(\gamma)$ in (\ref{pathlength}). And $\gamma=(\gamma_x,\gamma_{xy}^\circ,\gamma_y)$ is a minimising path for the geodetic distance whose part in $N(\mathcal{A})$ connects $\sigma(x)$ with $\sigma(y)$. It exists by the argument in the beginning of the proof. Let now $z\in X$ be a third point. Then 
\begin{align*}
d_\Lambda(x,y)&\le\ell_\Lambda(\gamma_x)+d_\Lambda(\sigma(x),\sigma(z))+d_\Lambda(z,y)+\ell_\Lambda(\gamma_y)
\\
&\le\ell_\Lambda(\gamma_x)+d_\Lambda(\sigma(x),\sigma(z))+2\ell(\gamma_z)+d_\Lambda(\sigma(z),\sigma(y))
\\
&=d_\Lambda(x,z)+d_\Lambda(z,y)\,,
\end{align*}
where $\sigma(z)$ is the root of the tree of balls containing $z$, and $\gamma_z$ is the geodesic between $\sigma(z)$ and $z$ in that tree.

\smallskip
The assertion now follows.
\end{proof}

\end{document}
